% TOP_PROBLEM: {"id": 3235, "title": "Circular colouring the orthogonality graph", "category_id": 3, "category": {"id": 3, "name": "graph_theory", "display_name": "Graph Theory", "description": "Problems involving graphs, networks, and their properties.", "slug": "graph-theory", "order_index": 3, "created_at": "2026-07-31T15:26:25.671Z"}, "status": "open", "problem_number": "OPG-2040", "set_id": 12}
% FIRST_POSTED: 2026-10-01
% ENTRY_KIND: statement_only
% UnsolvedMath ID 3235; source catalogue code OPG-2040
% !TeX program = lualatex
% Definition and sources only; this file is not an LLM solution attempt.
\documentclass[11pt,a4paper]{article}
\usepackage[margin=25mm]{geometry}
\usepackage{fontspec}
\setmainfont{lmroman10-regular.otf}[BoldFont=lmroman10-bold.otf,
 ItalicFont=lmroman10-italic.otf,BoldItalicFont=lmroman10-bolditalic.otf]
\usepackage{amsmath,amssymb,mathtools}
\usepackage{unicode-math}
\setmathfont{latinmodern-math.otf}
\AtBeginDocument{\let\setminus\smallsetminus}
\usepackage{xurl,hyperref}
\hypersetup{colorlinks=true,urlcolor=blue}
\setcounter{secnumdepth}{0}
\setlength{\parindent}{0pt}
\setlength{\parskip}{5pt}
\setlength{\emergencystretch}{3em}
\begin{document}
\section*{Circular colouring the orthogonality graph}\label{3235}
{\small UnsolvedMath ID 3235 · OPG-2040 · Graph Theory · OpenGarden}
\subsection{Definitions and mathematical statement}
Let $\mathcal O$ be the graph whose vertices are the one-dimensional linear subspaces of $\mathbb R^3$. Two distinct lines $L,M$ are adjacent when they are perpendicular, meaning $x\cdot y=0$ for nonzero representatives $x\in L$ and $y\in M$. Scaling representatives, including by negative numbers, does not change this condition. Consequently opposite directions on the unit sphere correspond to one vertex, rather than two.

For integers $p\ge2q\ge2$, a circular $(p,q)$-coloring of $\mathcal O$ is a map $c:V(\mathcal O)\to\{0,\ldots,p-1\}$ such that $q\le|c(L)-c(M)|\le p-q$ whenever $L\perp M$. Define $\chi_c(\mathcal O)$ as the infimum of all ratios $p/q$ for which such a coloring exists. A coloring is an arbitrary function; no continuity, measurability, or definability is imposed.

The question is whether $\chi_c(\mathcal O)=4$. In particular, since an ordinary proper four-coloring gives the upper bound four, the disputed direction asks whether every rational ratio $p/q<4$ is impossible. It would suffice to produce, for each such ratio, a finite configuration of lines admitting no $(p,q)$-coloring, by the compactness principle for finite-color constraints.

The target is circular chromatic number, which may be nonintegral, and not merely the ordinary chromatic number. Ruling out ordinary three-colorings alone does not establish the desired lower bound of four for circular colorings.
\subsection{Short English statement}
Is the circular chromatic number of the graph of perpendicular lines through the origin in three-dimensional space equal to four?
\subsection{Sources}
\begin{itemize}
\item[S1] ULAM.AI and the UnsolvedMath Contributors, supplied problems.json, record 3235 (OPG-2040), OpenGarden. Catalogue identity and source transcription; CC BY 4.0.
\url{https://huggingface.co/datasets/ulamai/UnsolvedMath}.
\item[S2] Open Problem Garden, Circular colouring the orthogonality graph. Original problem and discussion.
\url{http://www.openproblemgarden.org/op/circular_colouring_the_orthogonality_graph}.
\end{itemize}
\end{document}
